\subsection{Relative homotopies}

Let X, Y be topological spaces, let A be a subspace of X, and let B be a subspace of Y.

Let \(f:\mathrm{X}\to\mathrm{Y}\) be a continuous map such that \(f(\mathrm{A})\subset\mathrm{B}\). We say that f is a homotopy equivalence of the pair (X,A) onto the pair (Y,B)

if there exists a continuous map \(g\colon Y\to X\) such that \(g(B)\subset A\), a homotopy \(\sigma\colon X\times I\to X\), fixed on A, relating \(\text{Id}_X\) to \(g\circ f\), and a homotopy \(\tau\colon Y\times I\to Y\), fixed on B, relating \(\text{Id}_Y\) to \(f\circ g\).

Suppose these conditions are satisfied. Then:

\begin{enumerate}
\def\labelenumi{\alph{enumi})}
\item
  The map f is a homotopy equivalence from X to Y, the map g is a homotopy equivalence from Y to X, and these homotopy equivalences are inverse to each other up to homotopy.
\item
  The map g is then a homotopy equivalence of the pair (Y,B) onto the pair (X,A), called an inverse of f up to homotopy.
\item
  The maps f and g induce, by restriction to subspaces, homeomorphisms from A onto B that are inverse to each other.
\end{enumerate}

We say that the pairs (X,A) and (Y,B) are homotopy equivalent if there exists a homotopy equivalence of the pair (X,A) onto the pair (Y,B). The relation “X, Y are topological spaces, A is a subspace of X, B is a subspace of Y, and the pairs (X,A) and (Y,B) are homotopy equivalent” is an equivalence relation.

